% @@STATS_BEGIN@@ — auto-generated by habor-analyze, do not edit manually
% Auto-generated by habor-analyze — do not edit manually.
% Re-generate: uv run habor-analyze studies

% ── General ──
\newcommand{\NumIncludedBenchmarks}{53}
\newcommand{\NumHaborMixTasks}{160}
\newcommand{\NumModels}{8}
\newcommand{\NumAgents}{4}

% ── Within-family: model vs agent effect ──
\newcommand{\AnthropicModelEffectMean}{29.042}
\newcommand{\AnthropicAgentEffectMean}{9040229646.466}
\newcommand{\AnthropicModelAgentRatio}{0.0}
\newcommand{\AnthropicModelWins}{41}
\newcommand{\AnthropicTotalBenchmarks}{47}
\newcommand{\AnthropicModelWinPct}{87}
\newcommand{\GoogleModelEffectMean}{4.006}
\newcommand{\GoogleAgentEffectMean}{4.628}
\newcommand{\GoogleModelAgentRatio}{0.9}
\newcommand{\GoogleModelWins}{27}
\newcommand{\GoogleTotalBenchmarks}{50}
\newcommand{\GoogleModelWinPct}{54}
\newcommand{\OpenAIModelEffectMean}{35.452}
\newcommand{\OpenAIAgentEffectMean}{41.405}
\newcommand{\OpenAIModelAgentRatio}{0.9}
\newcommand{\OpenAIModelWins}{35}
\newcommand{\OpenAITotalBenchmarks}{49}
\newcommand{\OpenAIModelWinPct}{71}

% ── Adjusted effects ──
\newcommand{\ModelEffectSpan}{8013572.44}
\newcommand{\AgentEffectSpan}{3672886.76}
\newcommand{\TopModel}{claude-opus-4-6}
\newcommand{\TopModelEffect}{1669494.17}
\newcommand{\BottomModel}{claude-haiku-4-5-20251001}
\newcommand{\BottomModelEffect}{-6344078.28}
\newcommand{\TopAgent}{terminus-2}
\newcommand{\TopAgentEffect}{-0.37}
\newcommand{\BottomAgent}{claude-code}
\newcommand{\BottomAgentEffect}{-3672887.13}

% ── Agent lift vs terminus-2 ──
\newcommand{\geminicliMeanDelta}{-0.085}
\newcommand{\geminicliWinRate}{60}
\newcommand{\codexMeanDelta}{-0.103}
\newcommand{\codexWinRate}{72}
\newcommand{\claudecodeMeanDelta}{-5342380.699}
\newcommand{\claudecodeWinRate}{62}

% ── Terminus delta by model×agent (top pairs) ──
\newcommand{\TerminusgptFiveFourcodexDelta}{0.822}
\newcommand{\TerminusgptFiveFourcodexWinRate}{94}
\newcommand{\TerminusgptFiveminicodexDelta}{0.504}
\newcommand{\TerminusgptFiveminicodexWinRate}{83}
\newcommand{\TerminusclaudesonnetFourSixclaudecodeDelta}{0.103}
\newcommand{\TerminusclaudesonnetFourSixclaudecodeWinRate}{64}
\newcommand{\TerminusgeminiThreeOnepropreviewgeminicliDelta}{0.088}
\newcommand{\TerminusgeminiThreeOnepropreviewgeminicliWinRate}{62}
\newcommand{\TerminusclaudeopusFourSixclaudecodeDelta}{-0.005}
\newcommand{\TerminusclaudeopusFourSixclaudecodeWinRate}{42}
\newcommand{\TerminusgeminiThreeflashpreviewgeminicliDelta}{-0.257}
\newcommand{\TerminusgeminiThreeflashpreviewgeminicliWinRate}{58}
\newcommand{\TerminusgptFivenanocodexDelta}{-1.635}
\newcommand{\TerminusgptFivenanocodexWinRate}{38}
\newcommand{\TerminusclaudehaikuFourFiveTwoZeroTwoFiveOneZeroZeroOneclaudecodeDelta}{-16027142.196}
\newcommand{\TerminusclaudehaikuFourFiveTwoZeroTwoFiveOneZeroZeroOneclaudecodeWinRate}{81}

% ── Benchmark headroom ──
\newcommand{\NumSaturatedBenchmarks}{12}
\newcommand{\NumHardBenchmarks}{7}
\newcommand{\MeanBestScore}{53.722}

% ── Per-domain headroom stats ──
\newcommand{\NumAgentsToolsAndSystemsSaturated}{0}
\newcommand{\NumAgentsToolsAndSystemsHard}{1}
\newcommand{\NumAgentsToolsAndSystemsBenchmarks}{7}
\newcommand{\MeanAgentsToolsAndSystemsBestScore}{65}
\newcommand{\NumDataAndAnalyticsSaturated}{0}
\newcommand{\NumDataAndAnalyticsHard}{1}
\newcommand{\NumDataAndAnalyticsBenchmarks}{1}
\newcommand{\MeanDataAndAnalyticsBestScore}{38}
\newcommand{\NumKnowledgeAndLongContextSaturated}{2}
\newcommand{\NumKnowledgeAndLongContextHard}{0}
\newcommand{\NumKnowledgeAndLongContextBenchmarks}{4}
\newcommand{\MeanKnowledgeAndLongContextBestScore}{83}
\newcommand{\NumMathematicsAndReasoningSaturated}{3}
\newcommand{\NumMathematicsAndReasoningHard}{0}
\newcommand{\NumMathematicsAndReasoningBenchmarks}{6}
\newcommand{\MeanMathematicsAndReasoningBestScore}{92}
\newcommand{\NumMultimodalSaturated}{0}
\newcommand{\NumMultimodalHard}{0}
\newcommand{\NumMultimodalBenchmarks}{1}
\newcommand{\MeanMultimodalBestScore}{69}
\newcommand{\NumProfessionalDomainsSaturated}{0}
\newcommand{\NumProfessionalDomainsHard}{0}
\newcommand{\NumProfessionalDomainsBenchmarks}{5}
\newcommand{\MeanProfessionalDomainsBestScore}{75}
\newcommand{\NumSafetyAndSecuritySaturated}{1}
\newcommand{\NumSafetyAndSecurityHard}{0}
\newcommand{\NumSafetyAndSecurityBenchmarks}{1}
\newcommand{\MeanSafetyAndSecurityBestScore}{97}
\newcommand{\NumScientificResearchSaturated}{1}
\newcommand{\NumScientificResearchHard}{1}
\newcommand{\NumScientificResearchBenchmarks}{9}
\newcommand{\MeanScientificResearchBestScore}{29502}
\newcommand{\NumSoftwareEngineeringSaturated}{5}
\newcommand{\NumSoftwareEngineeringHard}{4}
\newcommand{\NumSoftwareEngineeringBenchmarks}{16}
\newcommand{\MeanSoftwareEngineeringBestScore}{73}
\newcommand{\MostHeadroomDomainOne}{Data \& Analytics}
\newcommand{\MostHeadroomDomainOneScore}{38}
\newcommand{\MostHeadroomDomainTwo}{Agents, Tools \& Systems}
\newcommand{\MostHeadroomDomainTwoScore}{65}
\newcommand{\MostHeadroomDomainThree}{Multimodal}
\newcommand{\MostHeadroomDomainThreeScore}{69}
% ── SE hardest benchmarks (sorted by score) ──
\newcommand{\SEHardBenchmarkOne}{AlgoTune}
\newcommand{\SEHardBenchmarkOneScore}{33}
\newcommand{\SEHardBenchmarkTwo}{BigCodeBench}
\newcommand{\SEHardBenchmarkTwoScore}{44}
\newcommand{\SEHardBenchmarkThree}{GSO}
\newcommand{\SEHardBenchmarkThreeScore}{44}
\newcommand{\SEHardBenchmarkFour}{FeatureBench}
\newcommand{\SEHardBenchmarkFourScore}{46}

% ── Effective dimensionality ──
\newcommand{\ParticipationRatio}{1.0}
\newcommand{\ComponentsForNinety}{1}
\newcommand{\ComponentsForNinetyFive}{1}
\newcommand{\TopOneVariance}{100}
\newcommand{\TopThreeVariance}{100}

% ── Domain correlation ──
\newcommand{\WithinDomainMedianCorr}{0.68}
\newcommand{\CrossDomainMedianCorr}{0.63}
\newcommand{\WithinDomainFracAboveSeven}{45}
\newcommand{\CrossDomainFracAboveSeven}{37}
\newcommand{\WithinAgentsToolsAndSystemsMedianCorr}{0.75}
\newcommand{\WithinKnowledgeAndLongContextMedianCorr}{0.62}
\newcommand{\WithinMathematicsAndReasoningMedianCorr}{0.59}
\newcommand{\WithinProfessionalDomainsMedianCorr}{0.72}
\newcommand{\WithinScientificResearchMedianCorr}{0.54}
\newcommand{\WithinSoftwareEngineeringMedianCorr}{0.71}

% ── Greedy benchmark selection ──
\newcommand{\GreedyBenchmarksBelowSeven}{12}
\newcommand{\GreedyFirstAboveSeven}{bigcodebench}
\newcommand{\GreedyFirstAboveSevenStep}{13}

% ── HaborMix difficulty composition ──
\newcommand{\HaborMixMediumTasks}{108}
\newcommand{\HaborMixEasyTasks}{38}
\newcommand{\HaborMixHardTasks}{13}
\newcommand{\HaborMixFrontierTasks}{1}

% ── AIME case study ──
\newcommand{\AimeGptFourCodexBase}{0.733}
\newcommand{\AimeGptFourCodexAgent}{0.987}
\newcommand{\AimeGptFourCodexDeltaPP}{25.3}
\newcommand{\AimeGptMiniCodexBase}{0.893}
\newcommand{\AimeGptMiniCodexAgent}{0.960}
\newcommand{\AimeGptMiniCodexDeltaPP}{6.7}
\newcommand{\AimeGptNanoCodexBase}{0.773}
\newcommand{\AimeGptNanoCodexAgent}{0.743}
\newcommand{\AimeGptNanoCodexDeltaPP}{-3.0}

% ── Figure paths (all under figs/main/quantitative/) ──
\newcommand{\FigWithinFamilySummary}{figs/main/quantitative/within\_family\_model\_vs\_agent\_summary}
\newcommand{\FigWithinFamilyDetail}{figs/main/quantitative/within\_family\_model\_vs\_agent\_detail}
\newcommand{\FigModelAdjustedEffects}{figs/main/quantitative/benchmark\_model\_adjusted\_effects}
\newcommand{\FigAgentLiftHeatmap}{figs/main/quantitative/benchmark\_agent\_lift\_heatmap}
\newcommand{\FigTerminusDeltaByModel}{figs/main/quantitative/terminus\_delta\_by\_model\_heatmap}
\newcommand{\FigHeadroomByDomain}{figs/main/quantitative/benchmark\_headroom\_by\_domain}
\newcommand{\FigHeadroomSummary}{figs/main/quantitative/benchmark\_headroom\_tier\_summary}
\newcommand{\FigSimilarityHeatmap}{figs/main/quantitative/benchmark\_similarity\_clustered\_heatmap}
\newcommand{\FigUniquenessVsCoverage}{figs/main/quantitative/benchmark\_uniqueness\_vs\_coverage}
\newcommand{\FigEffectiveDimensionality}{figs/main/quantitative/benchmark\_effective\_dimensionality}
\newcommand{\FigGreedySelection}{figs/main/quantitative/benchmark\_greedy\_selection}
\newcommand{\FigTaskSimilarityHeatmap}{figs/main/quantitative/task\_similarity\_benchmark\_pair\_heatmap}
\newcommand{\FigTaskPredictability}{figs/main/quantitative/task\_hard\_to\_predict\_ranked}
\newcommand{\FigTaskRepresentatives}{figs/main/quantitative/task\_best\_representatives}
\newcommand{\FigTaskDifficultyComposition}{figs/main/quantitative/task\_reliable\_difficulty\_composition}
\newcommand{\FigHaborMixSelection}{figs/main/quantitative/harbormix\_selection\_diagnostics}

% @@STATS_END@@

\subsubsection{Benchmarks}

%%% !todo: need to make sure all the figures and data being used is from the finalized benchmarks.

%%% for each part of these, it can be included in the appendix for all the math equations, data details, and more figures.

%%%%%%%%%%%%%%%%%%%%%%%%%%%%%%%%%%%%%%%%%%%%%%%%%%%%%%%%%%%%%%%%%%%%%%%%%%%%%%%
\paragraph{How fast are models improving over time?}
% TODO: progress-over-time figure (needs model release dates in data).
% \lin{We need one main figure here: progress by the years of some main
% benchmarks that we care and can make apple-to-apple comparison. Each domain
% should share the same color - different benchmarks should have different
% color intensity.}

%%%%%%%%%%%%%%%%%%%%%%%%%%%%%%%%%%%%%%%%%%%%%%%%%%%%%%%%%%%%%%%%%%%%%%%%%%%%%%%
\paragraph{Which benchmarks still challenge frontier systems?}
Of the \NumIncludedBenchmarks{} (\outlinecomment{xiangning: these need to use the finalized benchmarks})benchmarks,
\NumSaturatedBenchmarks{} have best-system scores above 95\%, concentrated in
Mathematics \& Reasoning (AIME, KUMO, IneqMath),
Software Engineering (CompileBench, HumanEvalFix,
CRUST-Bench, QuixBugs, USACO), and Knowledge \& Long
Context (GPQA Diamond, SimpleQA). These benchmarks offer
little room to differentiate frontier systems.

The domains with the most headroom are \MostHeadroomDomainOne{} (mean best
score \MostHeadroomDomainOneScore\%) and \MostHeadroomDomainTwo{}
(\MostHeadroomDomainTwoScore\%), where no benchmark is saturated.
Software Engineering is the most heterogeneous domain: it contains
\NumSoftwareEngineeringSaturated{} of the \NumSaturatedBenchmarks{} saturated
benchmarks and \NumSoftwareEngineeringHard{} of the \NumHardBenchmarks{} hardest
benchmarks (best score $<$50\%), including \SEHardBenchmarkOne{}
(\SEHardBenchmarkOneScore\%), \SEHardBenchmarkTwo{} (\SEHardBenchmarkTwoScore\%),
\SEHardBenchmarkThree{} (\SEHardBenchmarkThreeScore\%), and
\SEHardBenchmarkFour{} (\SEHardBenchmarkFourScore\%). This spread suggests that
software engineering tasks vary widely in the capabilities they require, from
near-solved code repair to open-ended algorithmic problem solving.

\begin{figure}[t]
  \centering
  \includegraphics[width=\linewidth]{\FigHeadroomByDomain}
  \caption{Best system score per benchmark, colored by domain. The shaded area
    represents headroom---room for improvement. Benchmarks are sorted within
    each domain by best score.}
  \label{fig:headroom-by-domain}
\end{figure}

%%%%%%%%%%%%%%%%%%%%%%%%%%%%%%%%%%%%%%%%%%%%%%%%%%%%%%%%%%%%%%%%%%%%%%%%%%%%%%%
\paragraph{How much overlap exists across benchmarks?}

Using a BenchPress-style predictability analysis~\cite{benchpress2025}
(Appendix~\ref{app:benchpress}), we find that the benchmark space is
dominated by a single general-capability axis: PC1 alone explains
\TopOneVariance\% of cross-system variance, and most benchmark scores
can be accurately reconstructed from a handful of others via logit-space
regression blended with rank-2 SVD completion.

Crucially, \emph{unpredictability does not align with difficulty}.
Nearly-solved benchmarks like StrongReject and CompileBench are among the
hardest to predict---their ranking patterns are orthogonal to the
general-ability axis---while hard unsolved benchmarks like AlgoTune
(67\% headroom) and HLE (46\% headroom) are trivially predictable from
other scores. This decoupling motivates HarborMix
(Section~\ref{sec:harbormix}): a good meta-benchmark must prioritize
informativeness (orthogonal signal) over raw difficulty.

%%%%%%%%%%%%%%%%%%%%%%%%%%%%%%%%%%%%%%%%%%%%%%%%%%%%%%%%%%%%%%%%%%%%%%%%%%%%%%%
\subsubsection{Agent Harness}
\label{sec:agent-harness}
\outlinecomment{0.5 pages}

\paragraph{Does the model or the agent matter more?}
Within each model provider family (Appendix~\ref{app:within-family}), switching
the underlying model produces an
average score change \AnthropicModelAgentRatio{}$\times$ larger than switching
the agent scaffold for Anthropic (\AnthropicModelWinPct\% of benchmarks
model-dominant), \OpenAIModelAgentRatio{}$\times$ for OpenAI
(\OpenAIModelWinPct\%), and \GoogleModelAgentRatio{}$\times$ for Google
(\GoogleModelWinPct\%). Across all \NumModels{} models, the adjusted model
effect (Appendix~\ref{app:adjusted-effects}) spans \ModelEffectSpan{}
robust-scaled units---from \TopModel{}
($+$\TopModelEffect{}) to \BottomModel{} (\BottomModelEffect{})---compared to
\AgentEffectSpan{} across \NumAgents{} agents, a ratio exceeding
10$\times$. Model selection is therefore the primary lever for performance on
the majority of benchmarks, though the magnitude of this asymmetry varies
substantially by provider.

%%%%%%%%%%%%%%%%%%%%%%%%%%%%%%%%%%%%%%%%%%%%%%%%%%%%%%%%%%%%%%%%%%%%%%%%%%%%%%%
\paragraph{Does the agent harness help, and when?}
We compare each agent scaffold against the Terminus~2 neutral baseline across
all benchmarks. The clearest evidence that agent scaffolding requires a
minimum model capability comes from holding the agent fixed (Codex) across
three OpenAI models of varying strength:

\begin{center}\small
\begin{tabular}{lccc}
\toprule
Model & Terminus mean & Agent $\Delta$ & Win rate \\
\midrule
GPT-5.4 (strong)   & 0.498 & $+$0.184 & 94\% \\
GPT-5-Mini (mid)   & 0.402 & $+$0.119 & 81\% \\
GPT-5-Nano (weak)  & 0.265 & $-$0.024 & 34\% \\
\bottomrule
\end{tabular}
\end{center}

\noindent The relationship is perfectly monotonic (Spearman $\rho = 1.0$
between model strength and agent delta). GPT-5-Nano paired with Codex
\emph{loses} to the neutral harness on 66\% of benchmarks---agent overhead
actively degrades a model that lacks the base capability to leverage it.
Conversely, GPT-5.4 wins on 94\% of benchmarks with a substantial
$+$18.4~pp mean lift across 52 benchmarks.

The same threshold effect appears for Claude~Code: it lifts Claude-Haiku
(81\% win rate, $+$11.2~pp) but shows diminishing returns on the already-strong
Claude-Opus (40\% win rate, $+$3.1~pp), where the orchestration overhead
begins to outweigh gains on benchmarks that Opus already handles well with the
simpler harness. Taken together, these results suggest that agent scaffolding is
a \emph{capability multiplier}, not a capability substitute---it amplifies
models above a competence threshold but introduces net harm below it.

% \lin{I don't think we actually need a figure here - but if we need, a thin
% one should be good to show how different harnesses has an impact on different
% benchmarks/models.}

% Optional thin figure: agent lift heatmap/Users/linxiangning/Desktop/projects/tb_adapter_project/habor-mix-analyzer/figs/main/quantitative/benchmark_agent_lift_heatmap.png
% \begin{figure}[t]
%   \centering
%   \includegraphics[width=\linewidth]{\FigAgentLiftHeatmap}
%   \caption{Agent lift vs.\ Terminus~2 by benchmark, grouped by domain. Green
%     indicates improvement; brown indicates degradation relative to Terminus~2.}
%   \label{fig:agent-lift-heatmap}
% \end{figure}
